\section{Keeping canonical cancellation inside an exact cell norm (NS-67)}
\label{sec:ns67-canonical-cells}

\paragraph{Status.}
The canonical growth criterion in Section~\ref{sec:ns66-canonical-growth}
is still open. Here the complete once-smoothed error is evaluated before
any absolute-value estimate is taken. This gives exact arithmetic cell
identities, a smaller elementary tail bound, and a finite growing-interval
version of the same growth target. It does not provide the missing
subpolynomial estimate. No finite table proves a cofinal bound.

Set $e_N=J(\chi+B_N)$ and $F_N(t)=e_N(1/t)$, $t>0$.
Write
\[
 b_N=Ns_N,\qquad c_{N,n}=\mu(n)-b_N\mathbf1_{n=N},\qquad
 q_N(j)=\sum_{\substack{n\mid j\\n\leq N}}\mu(n).
\]
The letter $b_N$ in this section denotes the canonical boundary
coefficient, not a target-load vector. The identity
$\sum_{n\leq N}c_{N,n}/n=0$ is exact.

\begin{proposition}[Complete physical cells]
\label{prop:ns67-cells}
We have $F_N(t)=0$ for $0<t\leq N$, and, for every integer $j\geq N$,
\begin{equation}
 F_N(t)=f_{N,j}+a_{N,j}\log(t/j),\quad j\leq t\leq j+1,
 \qquad
 a_{N,j}=1-\sum_{n\leq N}\mu(n)\left\lfloor\frac jn\right\rfloor
                   +b_N\left\lfloor\frac jN\right\rfloor.
 \label{eq:ns67-cell}
\end{equation}
The endpoint values are generated by
\[
 f_{N,N}=0,\qquad
 f_{N,j+1}=f_{N,j}+a_{N,j}\log(1+1/j).
\]
Equivalently, putting $d_N(j)=q_N(j)-b_N\mathbf1_{N\mid j}$,
\[
 a_{N,j}-a_{N,j-1}=-d_N(j),\qquad j>N.
\]
The initial slope is $a_{N,N}=b_N$, and therefore
\begin{equation}
 F_N(t)=b_N\log(t/N)
          -\sum_{N<m\leq t}d_N(m)\log(t/m),\qquad t\geq N.
 \label{eq:ns67-divisor-response}
\end{equation}
For $h_j=\log(1+1/j)$ and
$I_k(h)=\int_0^h u^ke^{-u}\,du$, the complete cell energy is
\begin{equation}
 C_{N,j}=\int_j^{j+1}\frac{F_N(t)^2}{t^2}\,dt
 =\frac{f_{N,j}^2 I_0(h_j)+2f_{N,j}a_{N,j}I_1(h_j)
                       +a_{N,j}^2 I_2(h_j)}j.
 \label{eq:ns67-cell-energy}
\end{equation}
In particular $D_{N,1}^2=\sum_{j=N}^\infty C_{N,j}$.
Every signed divisor contribution and every cross term remains present.
\end{proposition}
\begin{proof}
The support assertion is the exact exterior interpolation in NS-62.
For $t\geq1$ the complete expression is
\[
 F_N(t)=\log t+\sum_{n\leq N}c_{N,n}A(t/n),
 \qquad A(z)=z-k\log z+\log(k!),\quad k=\lfloor z\rfloor.
\]
On a unit cell all floors are constant. The coefficient of $t$ is
$\sum c_{N,n}/n=0$, and the coefficient of $\log t$ is precisely
$a_{N,j}$. Continuity of $A$ at integers gives the endpoint recurrence.
The floor jumps give the divisor expression; at the initial endpoint
$\sum_{n\leq N}\mu(n)\lfloor N/n\rfloor=1$, so $a_{N,N}=b_N$.
The change $t=je^u$ proves~\eqref{eq:ns67-cell-energy}.
\end{proof}

\begin{lemma}[An elementary complete Stirling remainder]
\label{lem:ns67-remainder}
For every real $z\geq1$,
\begin{equation}
 \left|A(z)-\tfrac12\log(2\pi z)\right|\leq\frac1{6z}.
 \label{eq:ns67-remainder}
\end{equation}
\end{lemma}
\begin{proof}
Let $R(z)=A(z)-\tfrac12\log(2\pi z)$ and
$\widetilde B_2(z)=\{z\}^2-\{z\}+1/6$.
Then $R'(z)=(\{z\}-1/2)/z$ almost everywhere.
Stirling's formula gives $R(z)\to0$ as $z\to\infty$.
Since $\widetilde B_2$ is continuous at integers and its derivative
is $2(\{z\}-1/2)$ off the integers, integration by parts gives
\[
 R(z)=\frac{\widetilde B_2(z)}{2z}
             -\frac12\int_z^\infty\frac{\widetilde B_2(u)}{u^2}\,du.
\]
Both terms are controlled by $|\widetilde B_2|\leq1/6$, proving
\eqref{eq:ns67-remainder}. There is no unestimated periodic tail.
\end{proof}

Define the exact signed leading coefficients and the absolute remainder
mass by
\[
 A_N=1+\tfrac12\sum_{n\leq N}c_{N,n}=1-\tfrac12N\eta_N,
 \qquad B_N^\flat=\tfrac12\sum_{n\leq N}c_{N,n}\log\frac{2\pi}{n},
 \qquad L_N=\sum_{n\leq N}n|c_{N,n}|.
\]
The superscript distinguishes $B_N^\flat$ from the canonical function $B_N$.
Then for $t\geq N$,
\begin{equation}
 F_N(t)=A_N\log t+B_N^\flat+R_N^\flat(t),\qquad
 |R_N^\flat(t)|\leq\frac{L_N}{6t}.
 \label{eq:ns67-full-tail}
\end{equation}
For $T\geq N$ put $v_N(T)=A_N\log T+B_N^\flat$ and
\[
 P_N(T)=\frac{v_N(T)^2+2A_Nv_N(T)+2A_N^2}{T},\qquad
 r_N(T)=\frac{L_N}{6\sqrt{3T^3}}.
\]
The entire tail obeys
\begin{equation}
 \left|\int_T^\infty\frac{F_N(t)^2}{t^2}\,dt-P_N(T)\right|
       \leq2\sqrt{P_N(T)}\,r_N(T)+r_N(T)^2.
 \label{eq:ns67-tail-energy}
\end{equation}
This follows by integrating the squared linear logarithm exactly and
applying Cauchy--Schwarz to its cross term with the bounded remainder.
The right side includes the remainder square; no cancellation with it
is presumed.

\begin{proposition}[The polynomial cutoff retains the growth target]
\label{prop:ns67-polynomial-cutoff}
Let
\[
 Q_N=\left(\sum_{j=N}^{N^2-1}C_{N,j}\right)^{1/2},\qquad N\geq2.
\]
Unconditionally,
\begin{equation}
 0\leq D_{N,1}-Q_N\ll\log(2N).
 \label{eq:ns67-cutoff-error}
\end{equation}
Consequently
\[
 \lim_{N\to\infty}\frac{\log(1+Q_N)}{\log N}=\beta_*-\frac12,
\]
and RH is equivalent to $Q_N=O_\varepsilon(N^\varepsilon)$ for every
$\varepsilon>0$. This is an exact finite growing-interval reformulation,
not a proven bound for $Q_N$.
\end{proposition}
\begin{proof}
The elementary divisor identity gives
\[
 b_N=1+\sum_{n\leq N}\mu(n)\{N/n\},\qquad |b_N|\leq N+1.
\]
Hence $\sum|c_{N,n}|\leq2N+1$, $|A_N|\leq N+3/2$,
$|B_N^\flat|\ll N\log(2N)$, and $L_N\ll N^2$.
At $T=N^2$, therefore, $P_N(T)^{1/2}=O(\log(2N))$ and
$r_N(T)=O(1/N)$. Equation~\eqref{eq:ns67-tail-energy} bounds the
complete tail norm by $O(\log(2N))$. Disjoint-support orthogonality
gives $D_{N,1}^2=Q_N^2+\text{tail energy}$ and thus
\eqref{eq:ns67-cutoff-error}.
Apply the exponent theorem~\ref{cor:ns66-growth-exponent}. When
$\beta_*=1/2$, $0\leq Q_N\leq D_{N,1}$ gives exponent zero.
When $\beta_*>1/2$, subtracting an $O(\log N)$ error does not change
any positive power lower bound coming from a zero to the right of the
line. Take the supremum over those zeros and combine with the upper
bound $Q_N\leq D_{N,1}$.
\end{proof}

\paragraph{Complete boundary/integral comparison.}
Write $w_N=N\eta_N$ and
$I_N=\int_1^N M(u)\rho_u\,du/u$.
Then the exact Abel identity is $e_N=\tau+I_N-w_N\sigma_N$.
Once the full energy and the mixed integrals have been enclosed,
\[
 \|\tau+I_N\|^2=D_{N,1}^2+2w_N\langle e_N,\sigma_N\rangle
                          +\frac{w_N^2}{N}\|\sigma_1\|^2,
\]
\[
 \|I_N\|^2=\|\tau+I_N\|^2+2
       -2\{\langle e_N,\tau\rangle+w_N\langle\sigma_N,\tau\rangle\}.
\]
These are comparisons of the complete terms, not of their absolute
integrands. The implementation integrates both mixed terms on every
cell and retains their complete tails. The identities make the
cancellation measurable without discarding it; they do not bound it
uniformly in $N$.

\begin{proposition}[The two complete Abel terms have a uniform conditioning bound]
\label{prop:ns67-abel-conditioning}
Put $c_\sigma=\|\sigma_1\|$ and $k_N=k_{N,1}$ from
Proposition~\ref{prop:ns62-smoothed-obstruction}. With
$U_N=\tau+I_N$ and $V_N=w_N\sigma_N$, we have
\begin{equation}
 \|V_N\|\leq c_\sigma D_{N,1}+\frac{2c_\sigma}{\sqrt N},\qquad
 D_{N,1}\leq\|U_N\|+\|V_N\|
 \leq(1+2c_\sigma)D_{N,1}+\frac{4c_\sigma}{\sqrt N}.
 \label{eq:ns67-conditioning}
\end{equation}
Moreover define $\widetilde V_N=(w_N-k_N)\sigma_N$ and
$\widetilde U_N=U_N-k_N\sigma_N$. Then
$e_N=\widetilde U_N-\widetilde V_N$ and
\begin{equation}
 \|\widetilde V_N\|\leq c_\sigma D_{N,1},\qquad
 D_{N,1}\leq\|\widetilde U_N\|+\|\widetilde V_N\|
                \leq(1+2c_\sigma)D_{N,1}.
 \label{eq:ns67-centered-conditioning}
\end{equation}
Thus cancellation between these two complete vectors cannot change the
norm-growth exponent. Their internal arithmetic cancellation remains
unestimated. This does not justify replacing the integral $I_N$ by
its absolute integrand.
\end{proposition}
\begin{proof}
The localized unitary witness already proved in NS-62 gives
$|w_N-k_N|/\sqrt N\leq D_{N,1}$.
The defining integral of $k_N$ has absolute value at most $2$, since
$|\sin(2\pi x)/(\pi x)|\leq2$ and
$N\int_0^{1/N}\log(1/(Nx))\,dx=1$.
Use $\|\sigma_N\|=c_\sigma/\sqrt N$ and the triangle inequality.
The centered inequalities follow with no additive error; both centered
vectors differ by exactly $e_N$, so no cross term is omitted.
\end{proof}

\paragraph{Complete finite checks.}
The physical cell implementation evaluates $N=8,16,64,512,4096$ at
256 and 384 bits, with the entire tail beyond $T=1024N$ enclosed by
\eqref{eq:ns67-tail-energy}. It also retains the mixed tails with
$\sigma_N$ and $\tau$. Independent complete Gram evaluations at
$N=8,16,64$ lie inside all three corresponding physical enclosures.
Doubling the physical cutoff at $N=16,512$ gives overlapping full
energies and mixed terms. The replay verifies $175$ cross-precision
quantities, $28$ cutoff comparisons and $9$ independent Gram comparisons.
In all five cases the centered ratio in
\eqref{eq:ns67-centered-conditioning} lies below the certified
constant $1+2c_\sigma<4.617$. These are complete finite norms;
no bounded cofinal subsequence is inferred.

\paragraph{What changed in the proposed cancellation route.}
The large finite ratios from separately bounding the two original Abel
vectors can include a harmless $N^{-1/2}$ component. Removing its explicit
coefficient $k_N$ makes their comparison uniformly controlled by the
actual error. Therefore the proposed search for a power-saving hidden
solely in cancellation between those two \emph{complete} vectors does
not supply a weaker arithmetic target. The possible useful cancellation
is inside the integral, or inside the exact divisor-cell response
\eqref{eq:ns67-divisor-response}. No bound for either signed quantity
is proved here. This is an obstruction to that specific decomposition
argument, not a failure of the canonical subpolynomial-growth target.

\section{A positive renewal equation retains the arithmetic forcing (NS-68)}
\label{sec:ns68-renewal}

\paragraph{Status.}
The first multiplicative interval of the canonical response gives a
particularly simple signed arithmetic quantity. Its exact renewal
operator is positive and contractive in every strictly positive power
weight. This controls inversion of the operator but does not improve
the growth exponent of its forcing. The required arithmetic estimate
is still RH-equivalent. The argument below is a direct calculation;
no claim of priority for these reformulations is made.

For real $x>0$ put
\[
 M(x)=\sum_{n\leq x}\mu(n),\quad s(x)=\sum_{n\leq x}\frac{\mu(n)}n,
 \quad p(x)=x s(x)-M(x),
\]
where all three are zero for $x<1$. Then $p$ is continuous, including
at integers, and $p(x)=x\eta(x)$ for $x\geq1$.
Let $v=\log2$ and define
\begin{equation}
 f(x)=x s(x)v+\sum_{x<n\leq2x}\mu(n)\log\frac{2x}{n}.
 \label{eq:ns68-forcing}
\end{equation}
In particular $f(x)=0$ for $x<1/2$.

\begin{proposition}[Exact first-interval response and renewal identity]
\label{prop:ns68-renewal}
For each integer $N\geq1$,
\begin{equation}
 f(N)=F_N(2N)=e_N(1/(2N)).
 \label{eq:ns68-canonical-sample}
\end{equation}
For every real $x>0$,
\begin{equation}
 f(x)=v p(x)+\int_x^{2x}\frac{M(u)}u\,du
     =p(2x)-(1-v)p(x)-\int_x^{2x}\frac{p(u)}u\,du.
 \label{eq:ns68-renewal-identity}
\end{equation}
The coefficients in the corresponding average have total mass
$(1-v)+v=1$, and both are nonnegative.
\end{proposition}
\begin{proof}
On $N<m<2N$, every proper divisor of $m$ is at most $N$, so
$q_N(m)=-\mu(m)$. The endpoint $m=2N$ has zero logarithmic weight.
Equation~\eqref{eq:ns67-divisor-response} therefore gives
\eqref{eq:ns68-canonical-sample}. It also holds for $N=1$ directly.
Partial summation gives
\[
 \sum_{x<n\leq2x}\mu(n)\log\frac{2x}{n}
       =\int_x^{2x}\frac{M(u)-M(x)}u\,du,
\]
including integer endpoints, since the coefficient vanishes at $2x$.
This proves the first identity. The continuous, piecewise differentiable
function $p$ satisfies $u p'(u)-p(u)=M(u)$ almost everywhere. Integrate
$p'(u)-p(u)/u$ from $x$ to $2x$ to get the second identity. These
identities also apply when the interval crosses $1$.
\end{proof}

\begin{proposition}[Weighted causal inversion, with no power gain]
\label{prop:ns68-weighted-inversion}
For $\theta>0$ define
\[
 q_\theta=(1-v)2^{-\theta}+\frac{1-2^{-\theta}}\theta<1.
\]
Write $P_\theta(t)=e^{-\theta t}p(e^t)$ and
$G_\theta(t)=e^{-\theta t}f(e^{t-v})$, both zero for $t<0$.
Then
\begin{equation}
 P_\theta(t)=(1-v)2^{-\theta}P_\theta(t-v)
       +\int_0^v e^{-\theta u}P_\theta(t-u)\,du+G_\theta(t).
 \label{eq:ns68-causal}
\end{equation}
For every $1\leq a\leq\infty$ and every $T>0$,
\begin{equation}
 (1-q_\theta)\|P_\theta\|_{L^a(0,T)}
 \leq\|G_\theta\|_{L^a(0,T)}
 \leq(1+q_\theta)\|P_\theta\|_{L^a(0,T)}.
 \label{eq:ns68-two-sided}
\end{equation}
These bounds also apply on $(0,\infty)$ whenever either side is
finite. In particular, a bound $|f(x)|\leq Cx^\theta$ on $x>0$
implies
\begin{equation}
 |p(x)|\leq\frac{2^{-\theta}C}{1-q_\theta}\,x^\theta.
 \label{eq:ns68-sup-transfer}
\end{equation}
No smaller power follows from this renewal identity for general forcing.
\end{proposition}
\begin{proof}
Substitute $x=e^{t-v}$ in~\eqref{eq:ns68-renewal-identity}. The
convolution kernel in~\eqref{eq:ns68-causal} is the positive measure
\[
 K_\theta=(1-v)2^{-\theta}\delta_v
              +e^{-\theta u}\mathbf1_{(0,v)}(u)\,du.
\]
Its mass is $q_\theta$. At $\theta=0$ the mass is one; for
$\theta>0$ all its mass is strictly damped at positive delays, so
$q_\theta<1$. Extend a restriction to $(0,T)$ by zero. Causality
makes its convolution agree with the original convolution on $(0,T)$.
Young's inequality and the triangle inequality give both sides of
\eqref{eq:ns68-two-sided}. The finite-interval estimate followed by
$T\to\infty$ justifies the assertion without assuming an unknown
infinite norm is finite. For the supremum bound,
$|G_\theta|\leq2^{-\theta}C$.

For an explicit scope check, take the nonarithmetic continuous function
$p_\theta(x)=(x^\theta-1)\mathbf1_{x\geq1}$ and define its forcing
by the renewal identity. For $x\geq1$ that forcing is exactly
\[
 f_\theta(x)=\left\{2^\theta-(1-v)
                  -\frac{2^\theta-1}{\theta}\right\}x^\theta
            =2^\theta(1-q_\theta)x^\theta.
\]
The coefficient is positive. Thus the inverse preserves this power,
which excludes a generic power improvement by positivity alone. This
is a test of the operator class, not a counterexample concerning the
actual M\"obius coefficients.
\end{proof}

\begin{proposition}[The single sample retains every nontrivial-zero obstruction]
\label{prop:ns68-scalar-criterion}
The following two statements are equivalent to RH:
\[
 \forall\varepsilon>0,\quad p(x)=O_\varepsilon(x^{1/2+\varepsilon}),
 \qquad
 \forall\varepsilon>0,\quad f(N)=O_\varepsilon(N^{1/2+\varepsilon}).
\]
The second assertion is only over integer $N$. No such bound is
proved here.
\end{proposition}
\begin{proof}
RH implies $M(x)=O_\delta(x^{1/2+\delta})$. The complete tail
formula for $\eta$ in NS-66 then bounds $p$ at the same power.
Conversely the Mellin identity
\begin{equation}
 \int_1^\infty p(x)x^{-s-1}\,dx
             =\frac1{(s-1)s\zeta(s)},\qquad \Re s>1,
 \label{eq:ns68-p-transform}
\end{equation}
continues analytically to $\Re s>1/2$ if all the stated power bounds
hold. A zero there would give a genuine pole, a contradiction.
The apparent singularity at $s=1$ is removable. The functional equation
then gives RH.

The renewal identity and~\eqref{eq:ns68-sup-transfer} transfer these
power bounds in both directions between $p$ and $f$ for real $x$.
For the discrete-to-continuous step, the exact divisor identity gives
$|s(x)|\leq1+1/x$ for $x\geq1$. Off the discrete breakpoints,
\[
 f'(x)=v s(x)+\frac{M(2x)-M(x)}x,
 \qquad |f'(x)|\leq2v+2<4.
\]
The sums defining $f$ match continuously at integers and half-integers:
a term entering at $2x=n$ has coefficient zero, and a term leaving
at $x=n$ is exactly replaced by the change in $xs(x)v$.
Therefore $f$ is globally $4$-Lipschitz on $[1,\infty)$.
Its integer values suffice for every displayed power bound. Compact
initial intervals change only the constant.
\end{proof}

For an independent analytic check, termwise integration on $\Re s>1$
gives the exact Mellin response
\begin{equation}
 \int_0^\infty f(x)x^{-s-1}\,dx
  =\frac{H(s)}{\zeta(s)},\qquad
 H(s)=\frac{v}{s-1}+\frac{2^s-1-sv}{s^2}.
 \label{eq:ns68-forcing-transform}
\end{equation}
The factor $H$ has no zero on $\Re s>0$ (and has a pole at $1$).
Indeed the unweighted probability kernel has Laplace transform
\[
 \widehat K(s)=(1-v)2^{-s}+\frac{1-2^{-s}}s,
 \qquad H(s)=\frac{2^s\{1-\widehat K(s)\}}{s(s-1)}.
\]
For $\sigma=\Re s>0$, positivity of that kernel gives
$|\widehat K(s)|\leq q_\sigma<1$. Hence the numerator never
vanishes in this half-plane. The apparent pole of $H/\zeta$ at
$s=1$ is removable, as in~\eqref{eq:ns68-p-transform}.
Thus no nontrivial-zeta-zero contribution is removed by this scalar
sampling kernel. This is the same weighted contraction in Mellin
coordinates, not an additional arithmetic estimate.

\begin{corollary}[Critical weighted total energy cannot be finite]
\label{cor:ns68-critical-energy}
Put $q=q_{1/2}$. For every $X\geq1$,
\begin{equation}
 (1-q)^2\int_1^X\eta(x)^2\,dx
 \leq\frac12\int_{1/2}^{X/2}\frac{f(x)^2}{x^2}\,dx
 \leq(1+q)^2\int_1^X\eta(x)^2\,dx.
 \label{eq:ns68-finite-energy-comparison}
\end{equation}
Both integrals diverge as $X\to\infty$, unconditionally. Finiteness
of this particular critical weighted total energy is therefore not a
possible sufficient estimate to pursue. This does not exclude a
logarithmic or subpolynomial growth bound for these integrals, and is
not a claim that the canonical norm $D_{N,1}$ diverges.
\end{corollary}
\begin{proof}
Use~\eqref{eq:ns68-two-sided} with $a=2$, $\theta=1/2$ and
$T=\log X$, and change variables. Here
$\|P_{1/2}\|^2=\int_1^X\eta(x)^2\,dx$ and
$\|G_{1/2}\|^2=\tfrac12\int_{1/2}^{X/2}f(x)^2x^{-2}\,dx$.

If $\int_1^\infty\eta(x)^2dx$ were finite, Cauchy--Schwarz would
make~\eqref{eq:ns68-p-transform} analytic on $\Re s>1/2$ and give,
for $s=1/2+\delta+i\gamma$,
\[
 \left|\int_1^\infty p(x)x^{-s-1}dx\right|
       \leq\frac{\|\eta\|_2}{\sqrt{2\delta}}.
\]
At any known critical-line zero $1/2+i\gamma$ its meromorphic
continuation has a pole of integer order at least one. Along the
horizontal approach that grows at least as $\delta^{-1}$, contradicting
the displayed bound. The positive lower factor in
\eqref{eq:ns68-finite-energy-comparison} transfers divergence to the
forcing energy. No assumption on simplicity or RH is needed.
\end{proof}

\paragraph{Finite verification.}
Exact Maxima checks retain integer and noninteger endpoints, the physical
canonical sample, the Mellin response and both energy primitives.
Arb replays at 256 and 384 bits verify the complete renewal identity
and finite energy comparison for
$N=16,64,512,4096,65536,1048576$, using every M\"obius coefficient
through $2N$. There is no omitted tail in these finite identities.
The critical kernel mass and supremum-transfer constant satisfy
\[
 .8027641<q_{1/2}<.8027642,\qquad
 3.58508<\frac{2^{-1/2}}{1-q_{1/2}}<3.58509.
\]
The archive matches $74$ cross-precision quantities. No growth bound
is inferred from these six values.

\paragraph{Stop condition.}
The averaging geometry yields an explicit inverse bound, not a new
arithmetic estimate. A square-root-plus-epsilon estimate for its signed
forcing is exactly strong enough for RH and is still absent. Bounding
all critical weighted energy by a finite constant is actually false.
The remaining permissible objectives are growth estimates with an
explicit rate; neither positivity of the renewal kernel nor the finite
checks supply them.

\section{The optimized endpoint remainder is not a small correction on the tested blocks (NS-69)}
\label{sec:ns69-actual-remainder}

\paragraph{Scope.}
The coefficients in this section minimize the actual smoothed NB error.
They are not canonical M\"obius coefficients. We retain the complete
Gram, the old-space projection and the signed arithmetic remainder.
The finite certificates show a limitation of the two-moment model;
they neither prove an asymptotic obstruction to that model nor supply
the missing cofinal correlation estimate.

Use $R_N=\tau-\sum_{m\leq N}c_m\sigma_m$, $E_N=\|R_N\|^2$,
and $g_n=\langle R_N,\epsilon_n\rangle$, $N<n\leq2N$.
The $c_m$ solve the complete normal equations. Write
\[
 L_{N,n}=\frac{a_N\{\tfrac12\Delta_n(\log^2)
                  +(2-\gamma)\Delta_n\log\}
                   +\ell_N\Delta_n\log}{n},
 \qquad g_n=L_{N,n}+r_{N,n},
\]
with the exact $a_N,\ell_N$ from
\eqref{eq:ns61-residual-endpoint}. Thus $L$ denotes the endpoint
model and $r$ its signed error, rather than an independent substitute
for the actual correlation vector.

\begin{proposition}[Complete signed remainder and improved absolute allowance]
\label{prop:ns69-remainder}
With the complete arithmetic function $S_2$ from Ehm's Gram formula,
\begin{equation}
 r_{N,n}=-\sum_{m\leq N}\frac{c_m}{m}
       \left\{S_2(n/m)-\frac{n-1}{n}S_2((n-1)/m)\right\}.
 \label{eq:ns69-signed-remainder}
\end{equation}
For every $n>N$,
\begin{equation}
 |r_{N,n}|\leq
 \frac{\pi^2\sum_{m\leq N}m|c_m|}{144n^2(n-1)}.
 \label{eq:ns69-remainder-bound}
\end{equation}
This improves the earlier elementary allowance by a factor of $18$.
It does not assert that the allowance is small relative to $L$ or $g$.
\end{proposition}
\begin{proof}
Subtract the two complete Gram columns at $n,n-1$ with weights
$1,-(n-1)/n$. Their $K_2$ terms cancel exactly. Combining their
remaining polynomial terms with the exact target-load difference
gives $L_{N,n}$, while their $S_2$ terms give
\eqref{eq:ns69-signed-remainder}. This retains every old coefficient.
Alternatively, Lemma~\ref{lem:ns67-remainder} gives the physical
error bound
$|R_N(x)-a_N\log(1/x)-\ell_N|\leq x\sum m|c_m|/6$
on the support of $\epsilon_n$. Its nonnegativity and exact moment
$\int x\epsilon_n(x)dx=\pi^2/[24n^2(n-1)]$ give
\eqref{eq:ns69-remainder-bound}. Both derivations apply to the
complete functions, with no finite physical cutoff.
\end{proof}

Let $H_N$ be the complete projected Gram of the new differences, so
the full improvement is $\Delta E_N=g^T H_N^{-1}g=E_N-E_{2N}$.
For any nonzero real column $v$, restricting the new correction to
that column gives the exact improvement
\begin{equation}
 Q_N(v)=\frac{(v^Tg)^2}{v^TH_Nv}\leq\Delta E_N.
 \label{eq:ns69-selected-gain}
\end{equation}
For a two-column matrix $W$ it gives
$(W^Tg)^T(W^TH_NW)^{-1}(W^Tg)$. We use the fixed columns
\[
 W_{n,0}=\frac{N^2\Delta_n\log}{n},\qquad
 W_{n,1}=\frac{N^2\{\tfrac12\Delta_n(\log^2)
                                      -\log N\Delta_n\log\}}{n}.
\]
Their span contains $L$ exactly. They are independent for $N\geq2$:
the ratio of the second to the first is
$\tfrac12\log(n(n-1))-\log N$, strictly increasing in $n$.
The old projection is subtracted before evaluating every quotient.

\paragraph{Certified finite comparisons.}
Arb calculations at 256 and 384 bits use the complete Bernoulli/Hurwitz
tail enclosure of Section~\ref{sec:ns61-gram-remainder}. All solves
retain ball errors. The table gives outward decimal intervals; the
archived JSON files retain the narrower balls. The final two columns
are fractions of the full improvement, not relative errors.
\begin{center}\small
\begin{tabular}{rcccc}
\toprule
$N$ & $\|r\|_2/\|L\|_2$ & $\cos(L,g)$ & two-column fraction
 & $Q_N(g)/\Delta E_N$\\
\midrule
16 & $(.6074,.6076)$ & $(.8477,.8479)$ & $(.1204,.1205)$ & $(.2332,.2333)$\\
32 & $(.9250,.9252)$ & $(.7982,.7983)$ & $(.2036,.2037)$ & $(.3258,.3259)$\\
64 & $(2.0203,2.0205)$ & $(.5966,.5968)$ & $(.0806,.0807)$ & $(.2723,.2724)$\\
128 & $(3.1748,3.1750)$ & $(.4592,.4594)$ & $(.0330,.0331)$ & $(.2946,.2948)$\\
\bottomrule
\end{tabular}\end{center}
Here $\|\cdot\|_2$ is the Euclidean norm of the finite correlation
vector, while $E_N$ is the physical Hilbert-space squared error.
The exact vector identity
$\|g\|_2^2=\|L\|_2^2+2L^Tr+\|r\|_2^2$ is checked;
no mixed term is omitted. At $N=64$ the optimized logarithmic
coefficient satisfies
$-.000554<a_N<-.000553$, whereas it is positive at the other
three displayed sizes. Positivity of this coefficient is therefore
not available for all optimized blocks.

The smaller absolute allowance is still loose: its vector norm divided
by the actual remainder norm lies between $216$ and $541$ on these
four blocks. But the failure to treat the remainder as small is not
only a failure of that allowance: the actual remainder norm exceeds
the model norm at $N=64,128$. Conversely, the mixed term is positive
at those two sizes, so this is not evidence of universal destructive
cancellation. At the isolated far index $n=16N$, the relative scalar
remainders are below $.017$ in all four cases. This finite far-field
observation provides no uniform or cofinal estimate.

\paragraph{What had to change.}
A lower bound obtained by treating the leading endpoint terms as the
whole residual is unsupported. Even allowing both endpoint directions
captures only about $3.3\%$ of the full block improvement at $N=128$.
The actual-correlation direction captures about $29.5\%$ there, and
must retain its arithmetic remainder. These are finite comparisons,
not a convergence theorem. The open input is still an independently
proved, nonsummable relative contraction along a cofinal sequence.

\section{The elementary Gram background requires a structural correction (NS-70)}
\label{sec:ns70-background}

\paragraph{Scope.}
The endpoint model in NS-69 comes from the elementary part of Ehm's
complete $q=2$ Gram formula. Its scalar correction is small in some
entries, but that does not make it a small quadratic-form perturbation.
We first identify and repair a mismatch in diagonal regularity. Even
after that repair, a uniform comparison on all coefficient vectors is
excluded. None of these statements is a lower bound or obstruction
for the actual optimized residual direction.

For positive real $u,v$ let $G(u,v)=\langle\sigma_u,\sigma_v\rangle$.
With $K_1,K_2$ from Section~\ref{sec:ns61-gram-remainder}, define
\[
 K(u,v)=\frac{K_2+K_1|\log(u/v)|+\tfrac14\log^2(u/v)}{\max(u,v)},
 \qquad S(u,v)=G(u,v)-K(u,v).
\]
At integer indices this is the full arithmetic remainder of the Ehm
formula. Set $d=K_2-2K_1$. The constants are certified below; in
particular $d>0$ and $K_1>3/2$.

\begin{proposition}[A small entrywise correction cancels the diagonal cusp]
\label{prop:ns70-cusp}
Both $G$ and $K$ are positive definite on every finite collection of
distinct positive dilations. For the adjacent coefficient vector
corresponding to
$\epsilon_n=\sigma_n-(n-1)\sigma_{n-1}/n$, write $G[\epsilon_n]$
and $K[\epsilon_n]$ for its two quadratic forms. Then
\begin{equation}
 G[\epsilon_n]\sim\frac\kappa{n^3},\qquad
 K[\epsilon_n]\sim\frac d{n^2},\qquad
 \frac{S[\epsilon_n]}{K[\epsilon_n]}\longrightarrow-1.
 \label{eq:ns70-cusp-asymptotic}
\end{equation}
Consequently no fixed $c>0$ can satisfy $G\geq cK$ on all finite
integer coefficient vectors. In particular, no relative remainder
bound $|S[a]|\leq\delta K[a]$ with a fixed $\delta<1$ can hold
on that class. This follows from local regularity and does not use
an assumption about the location of zeta zeros.
\end{proposition}
\begin{proof}
Write $r=\log(v/u)\geq0$. Normalizing by $\sqrt{uv}$ turns $K$
into the translation kernel
\[
 \phi(r)=e^{-|r|/2}(K_2+K_1|r|+r^2/4).
\]
Its Fourier density, with
$\phi(r)=(2\pi)^{-1}\int\widehat\phi(t)e^{itr}dt$, is
\begin{equation}
 w_K(t)=\frac{d t^4+(K_2/2-3/2)t^2
                         +K_2/16+K_1/8+1/8}{(t^2+1/4)^3}>0.
 \label{eq:ns70-background-symbol}
\end{equation}
This is obtained by differentiating
$\int e^{-a|r|}e^{-itr}dr=2a/(a^2+t^2)$ twice in $a$ and then
putting $a=1/2$. The three numerator coefficients are positive by
the recorded constant enclosures. Hence $K$ is strictly positive
definite; $G$ is already a Gram of independent dilations.

For $h_n=\log(n/(n-1))$ direct evaluation gives
\begin{equation}
 n^2K[\epsilon_n]
 =K_2-2K_1(n-1)h_n-\tfrac12(n-1)h_n^2\longrightarrow d.
 \label{eq:ns70-exact-adjacent-background}
\end{equation}
NS-61 gives $n^3G[\epsilon_n]\to\kappa$ and its complete upper
bound. Thus the ratio $G[\epsilon_n]/K[\epsilon_n]$ tends to zero,
proving all assertions. The cusp is explicit:
$\phi'(0^+)=K_1-K_2/2=-d/2$, while the actual smoothed Gram
has the additional order of cancellation in
\eqref{eq:ns70-cusp-asymptotic}.
\end{proof}

\paragraph{The cusp correction is exactly diagonal in adjacent differences.}
For $B(u,v)=1/\max(u,v)$, direct evaluation gives
\[
 B[\epsilon_m,\epsilon_n]=\frac{\mathbf1_{m=n}}{n^2},
 \qquad B[\sigma_j,\epsilon_n]=0\quad(j<n).
\]
Indeed for $j<n$ the right difference is
$1/n-((n-1)/n)/(n-1)=0$, including $j=n-1$.
The diagonal difference equals $1/n^2$. Consequently the repair
subtracts exactly $d\,\operatorname{diag}(n^{-2})$ from the raw
new-difference Gram and leaves every elementary old/new cross column
unchanged. In particular it cannot improve the endpoint model $L$
in NS-69: the correction cancels from those correlations identically.
Its role is to repair the energy assigned to correction directions.
The old model Gram also changes, so a projected model must still
recompute its complete Schur complement.

\begin{proposition}[Repairing the cusp still leaves the zeta-weight comparison]
\label{prop:ns70-repaired-background}
The corrected elementary kernel
\begin{equation}
 K_0(u,v)=K(u,v)-\frac d{\max(u,v)}
 =\frac{2K_1+K_1|\log(u/v)|+\tfrac14\log^2(u/v)}{\max(u,v)}
 \label{eq:ns70-repair}
\end{equation}
is strictly positive definite and has no first-derivative cusp.
It satisfies
\[
 K_0[\epsilon_n]\sim\frac{K_1-1/2}{n^3},\qquad
 \frac{G[\epsilon_n]}{K_0[\epsilon_n]}
                \longrightarrow\frac\kappa{K_1-1/2}.
\]
Nevertheless there are no constants $c>0$ or $C<\infty$ such that,
on all finite integer coefficient vectors, respectively
\[
 G\geq cK_0\qquad\text{or}\qquad G\leq CK_0.
\]
Both failures are unconditional. They exclude uniform all-coefficient
comparisons, not selected-direction estimates, finite-block inequalities
with size-dependent constants, or the optimized NB criterion.
\end{proposition}
\begin{proof}
The Fourier density of the repaired kernel is
\begin{equation}
 w_0(t)=\frac{(K_1-3/2)t^2+K_1/4+1/8}{(t^2+1/4)^3}>0.
 \label{eq:ns70-repaired-symbol}
\end{equation}
This proves strict positivity and the cancellation of the cusp.
In~\eqref{eq:ns70-exact-adjacent-background} replace $K_2$ by
$2K_1$ and use
$(n-1)h_n=1-1/(2(n-1))+O(n^{-2})$ and
$(n-1)h_n^2=1/(n-1)+O(n^{-2})$.
This gives the displayed adjacent asymptotic.

The complete Gram has Mellin density
\[
 w_G(t)=\frac{|\zeta(1/2+it)|^2}{(t^2+1/4)^2}.
\]
For clarity, a uniform quadratic-form inequality on all finite integer
coefficients would force the corresponding pointwise inequality between
$w_G$ and $w_0$. Here is the full localization argument. Both densities
are continuous, even and integrable. Their homogeneous kernels extend
to all positive real dilation parameters. Replace any finite set of real
parameters $u_j$ by integers nearest to $D u_j$ and multiply both forms
by $D$; kernel continuity lets $D\to\infty$. Thus an integer inequality
extends to all real parameters. Riemann sums then extend it to compactly
supported real coefficient densities in $r=\log u$. Choose long real
cosine packets in $r$. Their squared Fourier transforms form approximate
identities at the paired frequencies $\pm t_0$. Continuity near the
frequencies and integrability away from them give the pointwise inequality.
This reasoning retains all quadratic cross terms and uses arbitrary
finite real coefficients, not an arithmetic coefficient rule.

A known critical-line zero makes $w_G(t_0)=0$ while $w_0(t_0)>0$,
so the lower comparison is impossible. For the upper comparison,
\[
 \frac{w_G(t)}{w_0(t)}
 =|\zeta(1/2+it)|^2\,
   \frac{t^2+1/4}{(K_1-3/2)t^2+K_1/4+1/8}.
\]
The rational factor tends to the positive constant $(K_1-3/2)^{-1}$.
Known unbounded critical-line values, for example
\cite[Theorem 1]{Soundararajan2008}, therefore exclude any finite $C$.
The lower failure uses zeros on the line, not hypothetical off-line zeros.
Neither failure says that the actual Gram, or an original Weil form, has
negative energy.
\end{proof}

\paragraph{Complete constant and adjacent-vector certificates.}
Arb at 256 and 384 bits encloses every numerator coefficient in
\eqref{eq:ns70-background-symbol} and
\eqref{eq:ns70-repaired-symbol} strictly above zero, with
$.0091604<d<.0091605$. Complete two-vector Gram comparisons give
\begin{center}\small
\begin{tabular}{rcc}
\toprule
$n$ & $G[\epsilon_n]/K[\epsilon_n]$
    & $G[\epsilon_n]/K_0[\epsilon_n]$\\
\midrule
256 & $(.36207,.36209)$ & $(1.11185,1.11187)$\\
1024 & $(.11987,.11989)$ & $(1.11423,1.11425)$\\
4096 & $(.03260,.03262)$ & $(1.11498,1.11500)$\\
\bottomrule
\end{tabular}\end{center}
The true quadratic energy stays positive. The Bernoulli/Hurwitz
remainder is retained; a doubled series cutoff gives overlapping
values for all $34$ constants and quantities. The asymptotic assertions
follow from the proofs, not from this finite table.

\paragraph{What had to change.}
The elementary background had the wrong local regularity; removing
$d/\max(u,v)$ repairs that mismatch and recovers the correct adjacent
power $n^{-3}$. This repair is real but insufficient: the complete
arithmetic factor is still neither uniformly above nor uniformly below
the repaired background. A useful next bound must keep the particular
residual direction or supply an additional arithmetic estimate. The
finite endpoint remainder in NS-69 is not silently discarded.

\section{Using the repaired kernel to select a direction (NS-71)}
\label{sec:ns71-selected-preconditioner}

\paragraph{Question and scope.}
NS-70 excludes a uniform comparison over every coefficient vector.
It does not exclude using the repaired kernel to choose one direction,
then measuring that direction with the true complete Gram. The finite
prediction tested here is that such a choice captures more of the
optimized improvement than the unmodified correlation vector. A good
finite direction is a computational finding, not a cofinal estimate.

On the new block $N<n\leq2N$, let $D_N$ and $H_N$ denote the raw
and old-space-projected actual difference Grams. Let $D_{0,N}$ and
$H_{0,N}$ be the corresponding Grams for the repaired kernel $K_0$.
All four matrices are strictly positive definite. In defining $H_{0,N}$,
both the old projection and the new block use $K_0$ consistently;
this is not the actual Hilbert-space projection.
For the actual optimized correlation vector $g$, consider
\[
 v_{\rm raw}=D_{0,N}^{-1}g,\qquad
 v_{\rm projected}=H_{0,N}^{-1}g.
\]
Both are only rules for selecting coefficients. Every reported energy
uses $H_N$ or $D_N$, with the complete arithmetic remainder retained.

\begin{proposition}[Selected-direction improvement and the precise open input]
\label{prop:ns71-selected-input}
For either positive model matrix $A_N\in\{D_{0,N},H_{0,N}\}$ put
\[
 v=A_N^{-1}g,\quad m_N=g^TA_N^{-1}g,\quad
 C_N=\frac{v^TH_Nv}{v^TA_Nv},\quad
 C_N^{\rm raw}=\frac{v^TD_Nv}{v^TA_Nv}.
\]
Whenever $g\ne0$, all denominators are positive, and
\begin{equation}
 E_N-E_{2N}\geq Q_N(v)=\frac{m_N}{C_N}
                  \geq\frac{m_N}{C_N^{\rm raw}}>0.
 \label{eq:ns71-selected-ratio}
\end{equation}
If along $N_j=2^jN_0$ independently justified estimates give
\[
 \frac{m_{N_j}}{E_{N_j}}\geq\alpha_j>0,
 \qquad C_{N_j}\leq B_j,
 \qquad \sum_j\frac{\alpha_j}{B_j}=\infty,
\]
then $E_{N_j}\to0$ and RH follows by the fixed smoothed NB criterion.
The same statement holds using $C_{N_j}^{\rm raw}$ in place of $C_{N_j}$.
No such infinite sequence of estimates is proved here.
\end{proposition}
\begin{proof}
Since $A_Nv=g$, $v^Tg=v^TA_Nv=m_N$. Substitute in
\eqref{eq:ns69-selected-gain}. Orthogonal projection gives
$H_N\leq D_N$, proving the second inequality with all old-space
cross terms accounted for. The assumed bounds yield
$E_{2N_j}\leq(1-\alpha_j/B_j)E_{N_j}$.
The exact gain ensures $0<\alpha_j/B_j\leq1$; if equality occurs,
the error vanishes already. Otherwise the divergent sum makes the
product of contraction factors tend to zero. Monotonicity extends
convergence from the dyadic sequence to all $N$. The criterion in
Proposition~\ref{prop:ns61-rh-equivalence} then applies.
This uses a selected scalar ratio, not a uniform matrix comparison.
\end{proof}

For its geometric meaning, let $u=H_N^{-1/2}g$ and
$z=H_N^{1/2}v$. Then the exact fraction of the full improvement is
\[
 \frac{Q_N(v)}{E_N-E_{2N}}=
       \frac{(u^Tz)^2}{\|u\|_2^2\|z\|_2^2}.
\]
Thus the finite experiment measures an angle in the actual projected
energy, including the arithmetic terms. A Euclidean angle or a small
entrywise kernel remainder would not be a substitute.

\paragraph{Certified finite outcome.}
The table gives outward intervals for the fraction of the full
improvement $E_N-E_{2N}$, using the actual projected energy throughout.
\begin{center}\small
\begin{tabular}{rccc}
\toprule
$N$ & $v=g$ & $v=D_{0,N}^{-1}g$ & $v=H_{0,N}^{-1}g$\\
\midrule
16 & $(.2332,.2333)$ & $(.7167,.7169)$ & $(.8390,.8391)$\\
32 & $(.3258,.3259)$ & $(.8580,.8581)$ & $(.8252,.8254)$\\
64 & $(.2723,.2724)$ & $(.9593,.9594)$ & $(.9577,.9578)$\\
128 & $(.2946,.2948)$ & $(.9729,.9730)$ & $(.9781,.9782)$\\
256 & $(.3840,.3842)$ & $(.9597,.9598)$ & $(.9705,.9707)$\\
\bottomrule
\end{tabular}\end{center}
Both rules outperform $v=g$ at every checked size. The projected
rule captures more than $97\%$ of the full improvement at $N=128,256$;
neither efficiency is asserted monotone. At $N=256$ it reduces the
squared residual by a fraction in $(.08081,.08083)$, compared with
the full optimum in $(.08326,.08328)$. The corresponding numerator
and selected ratio in Proposition~\ref{prop:ns71-selected-input}
satisfy
\[
 .7668< m_{256}/E_{256}<.7669,\qquad
 9.4884<C_{256}<9.4886.
\]
These two finite bounds make explicit where the remaining arithmetic
problem sits; no size-independent bound is inferred.

All five sizes were evaluated at 256 and 384 bits with the full
Bernoulli/Hurwitz Gram tail. A doubled series cutoff at $N=256$
replays the largest problem. The verifier matches $115$ quantities
between precisions, $23$ at the doubled cutoff, and $16$ overlapping
quantities against the separate NS-69 implementation. The minimum
recorded accuracy is $88$ bits. All model and actual solves retain
their ball enclosures. The Maxima checks concern the exact quotient
and projection algebra, not the missing arithmetic estimates.

\paragraph{Stop condition.}
The model inversion is explicit and numerically useful on the checked
blocks. The arithmetic numerator $m_N/E_N$ and the selected true-energy
ratio still need cofinal estimates. Finite efficiency, by itself,
does not control either quantity as $N\to\infty$. NS-70's obstruction
to all-coefficient comparisons is respected throughout; no new
zero-free region, uniform Weil floor, G2 or RH proof follows.
